\begin{scanpage}{0178}{160}
diffraction ray concept is used for calculating the scattering of scalar and electromagnetic plane waves by a finite circular cone with a flat base and also by a cone having a spherical rounding off instead of a flat base. The resulting expressions are not applicable in the vicinity of certain irradiation and observation directions. In \S~17 we showed that the field scattered by a cone and by certain bodies of rotation is not expressed only in terms of the functions f and g, which refer to diffraction rays diverging from a wedge edge. This result evidently attests to the impossibility of complete calculation of the scattering characteristic with the diffraction ray concept.

Diffraction problems arising in antenna theory are usually distinguished by their great complexity, since the corresponding metal bodies (mirror, horn, etc.) have a complicated shape. Since the dimensions of these bodies and the dimensions of the radiating apertures are considerably larger than the wavelength, the application of physical diffraction theory to antenna problems is very promising. Only the first steps have been taken in this direction. Thus, Kinber [60, 61] performed a calculation of the decoupling and lateral radiation of mirror antennas. The feature specific to mirror antennas is that diffraction rays arising at the mirror's edge undergo multiple reflection on its concave surface. This multiple reflection was studied by Kinber in more detail as applied to the concave surface of a cylinder and sphere [62, 63].

Diffraction problems relating to an antenna dipole---a thin cylindrical conductor---are investigated in Chapter VII, and references to the literature are also given there.

In conclusion, let us say a few words about diffraction of short waves by smooth bodies. The basic principles relating to such problems were set forth in the fundamental works of Fok and Leontovich. These principles were established by the following methods of mathematical diffraction theory:
\end{scanpage}
